\chapter{\ifproject%
      \ifenglish Experimentation and Results\else การทดลองและผลลัพธ์\fi
  \else%
      \ifenglish System Evaluation\else การประเมินระบบ\fi
  \fi}

\section{การประเมินประสิทธิภาพของระบบ}

การประเมินประสิทธิภาพของระบบจะเป็นการทดสอบความสามารถของระบบโดยแบ่งออกเป็น 3 ด้าน ได้แก่ ระยะเวลาในการประมวลผลข้อมูล การประเมินประสิทธิภาพการทำงานของ Database และความถูกต้องในการประมวลผล ดังรายละเอียดต่อไปนี้

\subsection{ระยะเวลาในการประมวลผลข้อมูล}

การประเมินระยะเวลาในการประมวลผลข้อมูลจะเป็นการประเมินความรวดเร็วในการทำงานของระบบโดยการทดสอบจะแบ่งรูปแบบการวัดระยะเวลาออกเป็น 2 แบบ ได้แก่ การวัดระยะเวลาในแต่ละขั้นตอน และการวัดระยะเวลารวมของระบบ เพื่อให้สามารถเห็นทั้งขั้นตอนที่ใช้เวลาในการประมวลผลนานและเวลาโดยรวมที่ระบบใช้ โดยขั้นตอนที่นำมาพิจารณาประกอบด้วย

\begin{enumerate}        
    \item \textbf{การเตรียมภาพ (Image Preprocessing)} วัดระยะเวลาที่ใช้ในการปรับปรุงและเตรียมภาพก่อนเข้าสู่กระบวนการตรวจจับวัตถุ
    \item \textbf{การตรวจจับและแยกวัตถุ (Contour Detection)} วัดระยะเวลาที่ใช้ในการประมวลผลภาพเพื่อค้นหาขอบเขตของวัตถุและแยกสกรูแต่ละชิ้นออกจากพื้นหลัง รวมถึงการกำหนดกรอบล้อมรอบวัตถุ (Bounding Boxes)
    \item \textbf{การวัดคุณลักษณะและจำแนกประเภท} วัดระยะเวลาที่ใช้ในการคำนวณคุณลักษณะของสกรู และการจำแนกประเภทสกรูตามเงื่อนไขที่กำหนด
    \item \textbf{การนับจำนวนและสรุปผล} วัดระยะเวลาที่ใช้ในการนับจำนวนวัตถุแต่ละประเภทและจัดเตรียมผลการตรวจสอบสำหรับส่งกลับไปยัง Web Application
\end{enumerate}

และการระยะเวลารวมในการประมวลผล ก็จะเป็นการวัดระยะเวลาตั้งแต่ Web Application ได้รับภาพจนถึงกระบวนการประมวลผลเสร็จสิ้นและระบบส่งผลลัพธ์กลับไปยัง Web Application เพื่อแสดงผล

\subsection{การประเมินประสิทธิภาพการทำงานของ Database}

การประเมินประสิทธิภาพการทำงานของ Database จะครอบคลุมการตรวจสอบความถูกต้อง ความสอดคล้อง ความปลอดภัย และประสิทธิภาพการทำงาน โดยแบ่งการประเมินออกเป็น 3 ด้าน ดังนี้

\begin{enumerate}        
    \item \textbf{การทดสอบความถูกต้องและความสอดคล้องของข้อมูล (Data Integrity and Consistency Testing)} 
          \begin{itemize}        
              \item ตรวจสอบความถูกต้องของข้อมูลสินค้าเมื่อมีการเพิ่มและลดจำนวนสินค้าคงเหลือ
              \item ตรวจสอบความสัมพันธ์ระหว่างข้อมูลสินค้า และประวัติการดำเนินงาน รวมถึงผลลัพท์เมื่อมีการลบข้อมูลที่มีความเกี่ยวข้องกัน
              \item ตรวจสอบการป้องกันข้อมูลซ้ำซ้อนและข้อมูลที่ไม่ครบถ้วน
          \end{itemize}
    \item \textbf{การทดสอบความปลอดภัยของข้อมูล (Data Security Testing)}
          \begin{itemize}        
              \item ตรวจสอบการกำหนดสิทธิ์การเข้าถึงข้อมูลของผู้ใช้งานแต่ละประเภท
              \item ตรวจสอบการป้องกันการเข้าถึงและแก้ไขข้อมูลโดยไม่ได้รับอนุญาต
          \end{itemize}
    \item \textbf{การทดสอบประสิทธิภาพของฐานข้อมูล (Database Performance Testing)}
          \begin{itemize}        
              \item ตรวจสอบระยะเวลาในการเพิ่ม เรียกดู แก้ไข และลบข้อมูล รวมถึงการค้นหาและอัพเดตข้อมูลสินค้าคงเหลือ
              \item ตรวจสอบความสามารถในการรองรับคำขอจากผู้ใช้งานหลายคนในเวลาเดียวกัน
              \item ตรวจสอบว่าระบบสามารถรองรับการเพิ่มขึ้นของข้อมูลและจํานวนผู้ใช้งานได้โดยไม่กระทบต่อประสิทธิภาพ
          \end{itemize}
\end{enumerate}

\subsection{ความถูกต้องในการประมวลผล}

การประเมินความถูกต้องในการประมวลผลจะเป็นการทดสอบว่าระบบจะสามารถประมวลผลได้อย่างถูกต้องภายใต้การจัดวางสกรูที่แตกต่างกัน โดยการทดสอบจะดำเนินการกับรูปแบบการจัดวางสกรูที่แตกต่างกัน 4 กรณี ได้แก่

\begin{enumerate}        
    \item \textbf{กรณีมีสกรูเพียง 1 ชิ้น} ทดสอบความสามารถในการตรวจจับกรอบวัตถุ จำแนกประเภท และวัดขนาดของสกรูหนึ่งชิ้นในภาพ
    \item \textbf{กรณีมีสกรูหลายชิ้นประเภทเดียวกัน} ทดสอบความสามารถในการตรวจจับและนับจำนวนสกรูที่มีประเภทเดียวกัน โดยพิจารณาว่าระบบสามารถแยกวัตถุแต่ละชิ้นออกจากกันได้ครบถ้วนหรือไม่
    \item \textbf{กรณีมีสกรูหลายประเภทในภาพเดียวกัน} ทดสอบความสามารถในการตรวจจับและจำแนกสกรูหลาย Class ภายในภาพเดียวกัน รวมถึงตรวจสอบผลการนับจำนวนแยกตามประเภท
    \item \textbf{กรณีสกรูสัมผัสหรือซ้อนทับกัน} ทดสอบข้อจำกัดของกระบวนการตรวจจับและแยกวัตถุในกรณีที่ขอบเขตของสกรูอยู่ใกล้กันหรือทับซ้อนกัน เพื่อประเมินว่าปัญหาดังกล่าวส่งผลต่อการตรวจจับ การจำแนก และการนับจำนวนอย่างไร
\end{enumerate}